% !TeX root = ../main.tex

\sysusetup{
  keywords  = {细胞色素P450，酶底物特异性，EZSpecificity，图神经网络，泄漏控制外部评估},
  keywords* = {Cytochrome P450, Enzyme-substrate specificity, EZSpecificity, Graph neural network, Leakage-controlled external evaluation},
}

\begin{abstract}
  酶底物特异性预测是生物催化、合成生物学与药物研发中的核心任务。2025 年发表于 \textit{Nature} 的 EZSpecificity 模型是该任务的代表性先进方法，声称具有跨酶家族通用性；然而其训练数据集 ESIBank 在细胞色素 P450 超家族上覆盖稀缺（25{,}225 条酶中仅 389 个 P450，占 $1.5\%$），而 P450 在药物代谢、植物次生代谢与农药降解中应用价值重要，可作为检验该通用性的严格外部评估基准。本研究开展三方面工作：其一，针对现有 P450 数据库多数仅提供酶序列、或配体类型（底物/产物/抑制剂/辅因子）记录不明的普遍局限，通过文献检索、PDB 结构分析与多轮人工审核（辅以多智能体协作流程），对整合自 RCSB PDB、ESIBank、P450Rdb、PlantP450DB 与 PCPD 五个来源的每条酶--配体记录严格复核与分类，构建了目前规模最大、标注最详细的 P450 酶--底物特异性公开基准数据集，含 1{,}622 个酶、2{,}125 个底物、47{,}510 条可用样本（正 4{,}362 / 负 43{,}148），显著扩展植物 P450 数据覆盖；其二，以 InterPro 功能域识别 ESIBank 原始训练集中的 389 个 P450 并过滤测试集对应样本，建立泄漏控制的外部公平对比协议；其三，引入残基级几何向量感知（GVP）通路构成双图架构，以零初始化加延迟解冻保证基线等价。在 7{,}963 条过滤后测试集上，EZSpecificity 公开预训练模型 Test AUC $= 0.5586$、AUPR $= 0.1004$（接近 $1{:}10.2$ 正负比随机基线）；本研究领域专属基线达到 AUC $= 0.9205$、AUPR $= 0.6403$（$\Delta\mathrm{AUC} = +0.3619$）。双图架构在单随机切分下与基线相当（AUC $+0.0008$、AUPR $-0.0275$），未见显著收益。研究结果表明通用先进模型在训练集稀缺家族上存在明显的 未见酶的泛化瓶颈，领域专属数据重训可显著改善 P450 特异性预测表现。
\end{abstract}

\begin{abstract*}
  Enzyme--substrate specificity prediction is a core task in biocatalysis, synthetic biology, and drug discovery. EZSpecificity, a cross-attention graph neural network published in \textit{Nature} in 2025, is a leading state-of-the-art method that claims generality across diverse enzyme families. However, its training dataset ESIBank substantially under-represents the cytochrome P450 superfamily: only 389 P450 enzymes ($1.5\%$ of ESIBank's 25{,}225 total enzymes) are included, whereas P450s hold significant importance in drug metabolism, plant secondary metabolism, and pesticide degradation---making the family a rigorous external benchmark for testing the model's generality claim. This thesis pursues three lines of work. First, to address limitations of existing P450 databases---most provide only enzyme sequences or leave ligand types (substrate, product, inhibitor, cofactor) unspecified---each enzyme--ligand record integrated from RCSB PDB, ESIBank, P450Rdb, PlantP450DB, and PCPD was rigorously re-curated via literature retrieval, PDB structural analysis, and multi-round human review (assisted by a multi-agent collaboration pipeline). The resulting benchmark, to our knowledge currently the largest and most detailed public P450 enzyme--substrate specificity dataset, contains $1{,}622$ enzymes, $2{,}125$ substrates, and $47{,}510$ usable samples ($4{,}362$ positive / $43{,}148$ negative), with substantially expanded plant P450 coverage. Second, a leakage-controlled external evaluation protocol was established by strictly identifying the $389$ ESIBank-overlapping P450s through InterPro domain verification and filtering matching enzymes from the held-out test set. Third, a dual-graph architecture was designed on top of the original framework by introducing a residue-level Geometric Vector Perceptron (GVP) pathway, with a zero-initialization and delayed-unfreeze fusion design that preserves strict baseline equivalence at training step zero. On the $7{,}963$-sample filtered test set, the public EZSpecificity checkpoint achieved Test AUC $= 0.5586$ and AUPR $= 0.1004$ (close to the $\approx 0.09$ random baseline at $1{:}10.2$ ratio), whereas our domain-specific baseline reached Test AUC $= 0.9205$ and AUPR $= 0.6403$ ($\Delta\mathrm{AUC} = +0.3619$). The dual-graph model showed a marginal Test AUC improvement of $+0.0008$ but a Test AUPR drop of $-0.0275$ under single random split, with no significant benefit. These results suggest that general-purpose state-of-the-art models exhibit a held-out enzyme generalization bottleneck on training-underrepresented enzyme families, and that domain-specific retraining substantially improves P450 substrate specificity prediction.
\end{abstract*}

% 法语（第三种语言）摘要，本科不需要
%\begin{abstract**}
%\end{abstract**}
